% Source: Wald GR, physical PDF pp. 429--431
%         (printed pp. 416--418).
% Coverage: Section 14.4, Black Hole Thermodynamics;
%           equations (14.4.1)--(14.4.2).
% Figures/tables/footnotes: none.
% Status: source-reviewed and compile-clean.
% Modernization: all script symbols use \mathcal.
% Uncertainties: none.

\subsection{Black Hole Thermodynamics}\label{sec:14.4}

In section~12.5 we obtained a remarkable mathematical analogy between the ordinary
laws of thermodynamics and laws applying to black holes which were derived from
classical general relativity. As can be seen from Table~12.1, if one makes the formal
replacements \(E\longrightarrow M\), \(T\longrightarrow\alpha\kappa\), and
\(S\longrightarrow A/8\pi\alpha\) in the laws of thermodynamics, one obtains valid
laws applying to black holes. A hint that this relationship might be more than just a
formal analogy already arose from the fact that the analogous quantities \(E\) and
\(M\) actually represent the same physical quantity, namely, total energy. However,
since the thermodynamic temperature of a black hole in classical general relativity is
absolute zero, the physical analogy appeared to end there.

However, we now have seen that when quantum effects are taken into account, then in
a very real sense the thermodynamic temperature of a black hole is not zero but is
\(\kappa/2\pi\) in units where \(k=\hbar=G=c=1\). A black hole absorbs and
\lq\lq emits\rq\rq\ particles exactly like a perfect blackbody at that temperature.
Thus, \(T\) and \(\kappa/2\pi\) are not merely analogous quantities but again
represent the same physical quantity. Setting \(\alpha=(2\pi)^{-1}\), we see that
the remaining analogous quantities are \(S\) and \(\frac14A\). Does
\(\frac14A\) physically represent the entropy of a black hole?

We now shall argue that the answer to this question appears to be yes. The ordinary
second law of thermodynamics states that the total entropy of matter in the universe
never decreases. However, if a black hole is present, one would like to restrict
attention to matter outside black holes, since matter that falls in is \lq\lq swallowed
up\rq\rq\ by the singularity within the black hole and, in any case, it cannot be
measured by an external observer. However, one easily can make the total entropy,
\(S\), of matter outside black holes decrease by dropping matter into a black hole.
On the other hand, the area theorem of classical general relativity states that the
surface area, \(A\), of a black hole never decreases. However, on account of the
quantum particle creation process, this law can be violated since
\(\langle\widehat T_{ab}\rangle\) does not satisfy the energy condition assumed in
the proof of the area theorem. Indeed, as discussed in the previous section, an
isolated black hole eventually \lq\lq evaporates\rq\rq\ completely, thereby
decreasing its area to zero. Thus, when black holes are present and quantum effects
are taken into account, both the ordinary second law and the area theorem can be
violated. However, in the processes where \(\delta S<0\) due to loss of matter into
a black hole, we increase the black hole area, \(\delta A>0\). Similarly, in the
evaporation process where \(\delta A<0\), we increase the entropy of matter outside
black holes, \(\delta S>0\), by the emission of thermal radiation. Therefore, let us
define the \emph{generalized entropy}, \(S'\), by (Bekenstein 1973b, 1974)
\begin{equation}
  S'=S+\frac14 k\frac{c^3A}{G\hbar},
  \tag{14.4.1}\label{eq:14.4.1}
\end{equation}
where we have restored the constants \(k\), \(\hbar\), \(G\), and \(c\) in this
formula. The fact that a decrease in \(S\) seems always to be compensated by an
increase in \(A\) and, similarly, a decrease in \(A\) seems always to be compensated
by an increase in \(S\), suggests that in any process the \emph{generalized second
law},
\begin{equation}
  \delta S'\geq0,
  \tag{14.4.2}\label{eq:14.4.2}
\end{equation}
may be valid.

In fact, the law~\eqref{eq:14.4.2} was first proposed by Bekenstein (with an
undetermined numerical factor in the coefficient of \(A\) in
eq.~[\ref{eq:14.4.1}]) prior to the discovery of the quantum particle creation
effects. However, in the context of purely classical general relativity,
equation~\eqref{eq:14.4.2} could be violated by putting a black hole in a thermal
bath at a temperature lower than that formally assigned to the black hole, thereby
producing heat flow from a cold body (the bath) to a hotter body (the black hole).
Alternatively, one could put matter in a box and carefully lower the box to the
horizon of the black hole before dropping it in. By this latter procedure, the entropy
of matter inside the box still would be lost, but by \lq\lq redshifting away\rq\rq\
the energy in the box, the area increase of the black hole could be kept arbitrarily
small. However, when the quantum effects are taken into account, these methods for
violating the generalized second law do not work. In the first example, the black
hole will radiate more than it will absorb, so heat actually flows from the black hole
to the bath. In the second example, the effects of the radiation---described at the
end of the previous section---which is felt by any stationary body near the black hole
alters the transfer of energy into the black hole in just such a way as to ensure that
the area increase of the black hole always is large enough to keep
equation~\eqref{eq:14.4.2} satisfied (Unruh and Wald 1982). Thus, the generalized
second law appears to hold, at least insofar as it can be tested by
gedankenexperiments.

However, the generalized second law has a natural and simple interpretation. It can
be viewed as nothing more than the \emph{ordinary} second law of thermodynamics
applied to a system containing a black hole. In order to maintain this view, one must
take the final step of assigning \(\frac14A\) as the \emph{physical entropy} of a
black hole. If this final step is taken, the analogy between the laws of thermodynamics
and the laws of black hole physics no longer would be an analogy at all. The laws of
black hole thermodynamics (including the generalized second law) would be seen as
being nothing more than the \emph{ordinary} laws of thermodynamics applied to a
self-gravitating quantum system containing a black hole.

Thus, the apparent validity of the generalized second law strongly suggests that
\(\frac14A\) is the thermodynamic entropy of a black hole. However, the underlying
physical basis by which \(\frac14A\) arises as the black hole entropy remains
unclear. Since the entropy of an ordinary physical system is essentially the logarithm
of the number of microscopic states compatible with the observed macroscopic state,
the assignment of \(\frac14A\) as the black hole entropy seems to indicate that in a
full quantum theory of gravity the number of \lq\lq internal states\rq\rq\ of a black
hole will be \(N\sim e^{A/4}\). However, general arguments for the validity of the
second law of thermodynamics for ordinary systems are based on notions of the
\lq\lq fraction of time\rq\rq\ a system spends in a given macroscopic state. Since
the nature of time in general relativity is drastically different from that in
nongravitational physics, it is not clear precisely how the generalized second law
will arise even if \(\frac14A\) is a measure of the number of internal states of a
black hole.

Thus, we appear to be in a situation with regard to black hole thermodynamics which
is very similar to the situation with regard to ordinary thermodynamics prior to the
discovery of the underlying basis of these laws arising from statistical physics. We
have discovered the laws of black hole thermodynamics---in this case by calculations
and gedankenexperiments rather than by laboratory experiments---but the underlying
basis of these laws is not known and presumably will not be fully understood until we
have a quantum theory of gravitation. Nevertheless, the existence of the laws of black
hole thermodynamics indicates the likelihood of a deep connection between gravitation,
quantum theory, and statistical physics. It remains for future investigations to explore
this connection further.
